\subsection{CRITERIA OF SUBSTITUTION}

Formal mathematics contains only explicitly written assemblies. Nevertheless, even with the use of abbreviating symbols, the development of mathematics strictly in accordance with this principle would lead to extremely long chains of reasoning. For this reason we shall establish criteria relating to indeterminate assemblies; each of these criteria will describe once for all the final result of a definite sequence of manipulations on these assemblies. These criteria are therefore not indispensable to the theory; their justification belongs to metamathematics.

The development of metamathematics itself requires, in practice, the use of abbreviating symbols, some of which have already been indicated. Most of these symbols are also used in mathematics.

We shall make use of the following criteria, called the criteria of substitution:

CS1. Let A and B be assemblies and let x and \(x'\) be letters. If \(x'\) does not appear in A, then \((B|x)A\) is identical with \((B|x')(x'|x)A\) .

CS2. Let \(A, B\) , and \(C\) be assemblies and let \(x\) and \(y\) be distinct letters\footnote{In accordance with what was said in no. 1, the phrase ``x and y are distinct letters'' is an abuse of language: it means that x and y denote distinct letters in the assemblies under consideration.}. If \(y\) does not appear in \(B\) , then \((B|x)(C|y)A\) is identical with

\[
\left(\boldsymbol {C} ^ {\prime} | \boldsymbol {y}\right) (\boldsymbol {B} | \boldsymbol {x}) \boldsymbol {A},
\]

where \(C'\) is the assembly \((B|x)C\) .

CS3. Let A be an assembly and let x and \(x'\) be letters. If \(x'\) does not appear in A, then \(\tau_{\mathbf{x}}(A)\) is identical with \(\tau_{\mathbf{x}'}(A')\) , where \(A'\) is the assembly \((x'|x)A\) .

CS4. Let A and B be assemblies and let x and y be distinct letters. If x does not appear in B, then \((B|y)\tau_{x}A\) is identical with \(\tau_{x}(A')\) , where \(A'\) is the assembly \((B|y)A\) .

CS5. Let \(A, B, C\) be assemblies and let \(x\) be a letter. The assemblies \((C|x) (\neg A), (C|x) (\vee AB), (C|x) (\Longrightarrow AB), (C|x) (sAB)\) (where \(s\) is a specific sign) are respectively identical with \(\neg A', \vee A'B', \Longrightarrow A'B', sA'B',\) where \(A', B'\) are respectively \((C|x) A, (C|x) B\) .

As an example, let us indicate the principle of the verification of CS2. Compare the operation which takes us from A to \((B|x)(C|y)A\) with the operation which takes us from A to \((C'|y)(B|x)A\) . In each operation, no sign which appears in A and is distinct from x and y is altered. At every place where x appears in A, we have to substitute B for x in the first and in the second operation; this is clear in regard to the first operation, and in regard to the second it follows from the fact that y does not appear in B. Finally, at every place where y appears in A, the first operation consists in replacing C for y, then B for x at every place where x appears in C; but it is clear that this comes to the same thing as substituting for y, wherever it appears in A, the assembly \((B|x)C\) .

